\section{Conclusion}
\label{sec:conc}

\noindent In this paper, we study how the 1933 abrogation of gold clauses---and the ensuing leverage uncertainty that persisted until the Supreme Court's 1935 decision---affected real investment in 1933 and 1934. Firms with higher pre-episode gold clause exposure cut investment significantly in 1933 and 1934 relative to other firms, and the gap disappears in 1935 once the uncertainty is resolved. These patterns are consistent with the debt-overhang mechanism: legal uncertainty about the fate of the gold clause increased debtholders' claim in expectation, discouraging investment.

Firms exposed to gold clauses also increased equity payouts, driven by cash dividends, in 1933--1934 and reverted afterward, while profitability did not deteriorate differentially. We find some evidence that the impact of the abrogation was strongest for lower-rated issuers, consistent with higher default risk amplifying the wedge. Results are robust to industry--year fixed effects and rich controls that interact pre-1933 characteristics with year, and they do not arise when gold clause exposure through bond financing is replaced with other long-term debt components: exposure to bank debt and preferred equity does not generate a similar pattern. Excluding firms with bonds maturing in 1931--1934 yields similar estimates, indicating that the effect is not driven by refinancing pressures or contemporaneous bond market dislocations.

Taken together, our evidence shows that policy-driven legal uncertainty on the liability side can depress investment at scale, even in the absence of realized changes in leverage. By isolating a plausibly exogenous episode that shifts expected liabilities, we provide evidence consistent with dynamic debt overhang and quantify its aggregate relevance.

While this paper focuses on the impact of gold clause abrogation on large industrial firms, examining its effects on other sectors of the economy represents a promising avenue for future research. Evidence suggests that gold clauses were prevalent in farm mortgages, making the agricultural sector a particularly important case for investigation. Given that agriculture employed nearly a quarter of the workforce at the onset of the Depression, understanding how the abrogation affected farmers' debt burdens, investment decisions, and recovery patterns could provide crucial insights into the policy's broader economic impact. Similarly, public utilities and railroads---major issuers of gold-denominated corporate bonds and critical components of economic infrastructure---warrant further attention. Investigating how the abrogation influenced these sectors' financing costs, capital investment, and operational decisions would enhance our understanding of the full macroeconomic consequences of this landmark policy intervention and its role in shaping the trajectory of economic recovery during the 1930s.